\chapter{Introduction}
Modern deep learning relies almost exclusively on gradient-based optimization, with gradients computed end-to-end via backpropagation \cite{rumelhart1986learning}. However, this paradigm imposes two structural assumptions: every component of the computational graph must admit a gradient (or a tractable subgradient), and parameter updates must propagate through the chain rule. Beyond these technical constraints lies a more systemic concern: the overwhelming focus on gradient-based training risks a form of paradigm lock-in, whereby the community may overlook alternative learning frameworks that given comparable research investment could ultimately prove competitive or even superior. Many alternative paradigms have been proposed in the literature, ranging from classical proposals such as target propagation \cite{lee2015difference}, the method of auxiliary coordinates \cite{carreira2014distributed} and feedback alignment~\cite{lillicrap2016random}, to more recent schemes like equilibrium propagation \cite{scellier2017equilibrium}, forward-forward training \cite{hinton2022forward} and neuroevolution \cite{such2017deep}. Nevertheless, none of these paradigms have achieved widespread adoption, typically bottlenecked by their heuristic approximations and inability to computationally scale on modern hardware.

A promising, mathematically grounded alternative that has recently gained attention is to recast learning as a feasibility problem \cite{elser2019learning, bergmeister2025projection, elser2026learning}. Instead of minimizing a global loss function $\mathcal{L}$, this approach searches for some point $\tilde{\theta}$ in the intersection of multiple constraint sets. These sets are chosen in a way so that from a feasible point $\tilde{\theta}$, one can derive the parameters $\theta$ of the neural network such that the training data is perfectly interpolated. Conceptually, this paradigm shift takes the form
\begin{tcolorbox}[shift_card]
\[
\minimize_{\theta \in \Theta}\, \mathcal{L}(f(x;\theta), y)
\quad \Longrightarrow \quad
\textnormal{find } \tilde{\theta} \in \bigcap_{i=1}^{m} C_i.
\]\end{tcolorbox}
A rich literature exists to tackle such problems, where projection-based methods are the most well-known. Examples of such methods include, but are not limited to, cyclic projections \cite{gubin1967method}, Bregman projections \cite{bregman1967relaxation}, Dykstra's algorithm \cite{dykstra1983algorithm}, and averaged projections \cite{auslender_thesis}. In the two-set case, Douglas--Rachford splitting \cite{douglas1956numerical,eckstein1992douglas} has been studied extensively and applied across a wide range of problems \cite{bauschke2021douglasrachfordalgorithmsolvingpossibly, lindstrom2020surveyyearsdouglasrachford}. Convergence theory for these projection-based methods has been primarily developed in the convex setting \cite{bauschke1996projection}, and there have been recent advances in the nonconvex setting as well \cite{themelis2020douglas,attouch2013convergence,li2016douglas}.

While these projection-based algorithms are promising, they are currently not competitive with gradient-based methods even for shallow networks. Moreover they come at the cost of substantial memory overhead, which prohibits scaling at the hardware level. Further, the frameworks in \cite{elser2019learning,bergmeister2025projection} were developed largely in isolation from modern deep learning toolchains and therefore cannot benefit from the many software optimizations that underpin state-of-the-art libraries such as Torch~\cite{paszke2019pytorchimperativestylehighperformance}. In particular, while PJAX \cite{bergmeister2025projection} is built on top of JAX \cite{bradbury2018jax}, it utilizes only the high-performance tensor computation library while all neural network training logic has been completely redefined within PJAX itself. These limitations naturally lead us to our main research question.

%Concretely, there are currently three major limitations. First, the consensus step, explained in \cref{sec:nn_feas}, requires a duplication of the neural network parameters across all samples in a batch, yielding a severe memory consumption. Second, the graph-propagation scheme of \cite{bergmeister2025projection} requires splitting the computational graph into a bipartite graph for every batch, introducing a per-batch overhead absent from standard training pipelines. Third, the frameworks in \cite{elser2019learning,bergmeister2025projection} were developed largely in isolation from modern deep learning toolchains and therefore cannot benefit from the many software optimizations that underpin state-of-the-art libraries such as Torch \cite{paszke2019pytorchimperativestylehighperformance}. In particular, while PJAX \cite{bergmeister2025projection} is built on JAX \cite{bradbury2018jax}, it utilizes only the high-performance tensor computation library while all neural network training logic has been completely redefined within PJAX itself. These identified limitations naturally lead us to our research question.

%Recent work by Bergmeister \cite{bergmeister2025projection} and Elser \cite{elser2019learning} has shown that projection-based algorithms can train neural networks in practice, but at the cost of substantial memory overhead and with empirical performance that remains well below that of gradient-based training. A closer examination reveals three recurring bottlenecks. First, memory consumption is severe: the consensus step in the linear layer scales with $O(N_{\mathrm{in}}N_{\mathrm{out}}B)$ rendering training infeasible on commonly available hardware. Second, the graph-propagation scheme of~\cite{bergmeister2025projection} requires splitting the computational graph into a bipartite graph for every batch, introducing a per-batch overhead absent from standard training pipelines. Third, convergence is noticeably slower than Adam, even on small MLPs. Moreover, these frameworks were developed largely in isolation from modern deep learning toolchains and therefore cannot benefit from the many software optimizations that underpin state-of-the-art libraries such as Torch \cite{paszke2019pytorchimperativestylehighperformance}.\footnote{While PJAX is built on JAX, it utilizes the framework solely as a high-performance tensor computation library; all neural network training logic has been completely redefined within PJAX itself}.

\vspace{0.5em}
\begin{tcolorbox}[researchcard]
Can we design a projection-based training framework native to modern deep learning toolchains that eliminates severe memory bottlenecks, avoids computational overhead, and achieves performance competitive with gradient-based methods?
\end{tcolorbox}
\vspace{0.5em}

%In response to these challenges, we introduce $\mathcal{P}$Torch. The compatibility problem is resolved by building on PyTorch, and by performing graph traversal directly through the native \texttt{autograd}. Further, we introduce and motivate various approaches to improve the performance of projection-based method. As a result, the gap between gradient-based and projection-based optimization in our framework reduces, as closely as possible, to the gap between backpropagation and its projection-based counterpart isolating the optimization strategy from implementation artifacts. 


%Making projection-based learning a mainstream paradigm remains a long-term goal; as with gradient-based 
%optimization, the right implementation tricks to train deep networks such as weight-initialization like He initialization \cite{delving} for gradients have yet to be fully worked out.

\paragraph{Contributions} We address these challenges by introducing $\mathcal{P}$Torch. Concretely, our contributions are as follows.
\begin{enumerate}
    \item \textbf{A PyTorch-native projection framework.}
    We establish a conceptual link between cyclic projections, existing projection-based learning algorithms \cite{elser2019learning,bergmeister2025projection}, and computational graph traversal via PyTorch's \texttt{autograd}. Leveraging this connection, we develop a highly extensible, PyTorch-native projection framework that inherits PyTorch's software optimizations and enables training hybrid networks that benefit from the properties of both paradigms.
    
    \item \textbf{Eliminating memory overhead and improving performance.} To mitigate memory constraints, we introduce a fused projection strategy that avoids parameter duplications of prior methods \cite{elser2019learning, bergmeister2025projection}. Furthermore, we demonstrate that modern first-order optimizers, such as Adam \cite{kingma2014adam} and Muon \cite{jordan2024muon}, can be utilized as nonlinear relaxation schemes for cyclic projections. This insight yields substantial empirical gains, enabling our framework to match or exceed the performance of gradient-based baselines on shallow networks. 
    \item \textbf{Vanishing targets.}
    We prove that the base cyclic projection algorithm suffers from a problem analogous to vanishing gradients, thereby limiting its algorithmic scalability to deeper networks. Successfully scaling these models presents a clear avenue of future work: discovering the projection-based equivalents of innovative heuristics like He initialization \cite{delving}, batch normalization \cite{ioffe2015batch} and architectural modifications like skip connections \cite{he2016deep}.
\end{enumerate}
\section{Related work}
Besides projection-based learning, a large body of work has explored alternative gradient-free paradigms. We briefly discuss the most prominent ones here.
\subsection{Target propagation.}
Target propagation is a  gradient-free learning paradigm where layer-wise targets are propagated backward through the network instead of gradients \cite{lee2015difference}. While these methods attempt to generate these targets by approximating the inverse of the forward map, our projection-based approach instead treats layer-wise updates as local feasibility problems. Instead of relying on approximate inverses, $\mathcal{P}$Torch computes targets by mathematically projecting the current state onto the layer's constraint set to find the smallest adjustment required to satisfy the network's architectural constraints. Despite these different mechanisms for generating local targets, both approaches share a fundamental limitation: the attenuation of the learning signal with depth.  

\subsection{ADMM-based deep learning.}
The alternating direction method of multipliers (ADMM) \cite{neal2011distributed} offers a prominent alternative to the projection-based training of neural networks. Similar to projection-based methods, ADMM decouples the network into per-layer subproblems by introducing auxiliary variables and enforcing consistency constraints via a Lagrangian penalty. The dlADMM \cite{10.1145/3292500.3330936} algorithm updates parameters in a backward-then-forward sweep to propagate information more efficiently. Crucially, this approach reduces the subproblem complexity from $O(n^3)$ to $O(n^2)$ using iterative quadratic approximations and backtracking, while also establishing rigorous global convergence guarantees under mild conditions on the loss and activation functions. Ultimately, the layer-wise decomposition in dlADMM conceptually aligns with $\mathcal{P}$Torch's target propagation: both methods bypass end-to-end differentiation by decomposing the global objective into isolated, layer-wise optimization steps propagated through the computational graph.

\subsection{Gradient-free optimization frameworks.}
A variety of other gradient-free optimization methods have been proposed in the literature, including the method of auxiliary coordinates \cite{carreira2014distributed}. By unfolding the network into layers, this approach replaces a complex nested optimization problem with local optimization problems at the level of individual layers. Distinct paradigms such as neuroevolution have also attracted significant attention and even inspired dedicated frameworks \cite{tang2022evojax}. Another notable alternative is the more recent forward-forward algorithm \cite{hinton2022forward}. However, a common thread across these approaches is that none have yet achieved widespread adoption in mainstream deep learning or been able to train modern deep networks as efficiently as gradient-based methods.